\chapter{Introduction}
\label{chap:introduction}

Ceci est le premier chapitre de votre rapport. Vous pouvez utiliser du faux texte pour voir le rendu :

\lipsum[1-2] % Génère 2 paragraphes de faux texte

Vous pouvez référencer une figure (comme la \cref{fig:exemple}) ou une table (voir \cref{tab:exemple}).

\begin{figure}[h!] % [h!] = Tente de placer la figure "ici"
    \centering
    \includegraphics[width=0.6\textwidth]{logos/Logo_LICIT-ECO7_vertical.png}
    \caption{Ceci est un exemple de figure.}
    \label{fig:exemple}
\end{figure}

\begin{table}[h!]
    \centering
    \caption{Ceci est un exemple de tableau.}
    \label{tab:exemple}
    \begin{tabular}{lcc}
        \toprule
        Paramètre & Valeur & Unité \\
        \midrule
        Vitesse   & 100    & m/s   \\
        Pression  & 50     & Pa    \\
        \bottomrule
    \end{tabular}
\end{table}

\lipsum[3]